\chapter{KẾT LUẬN}
\label{chuong6}


Luận văn này đặt ra để trả lời một câu hỏi duy nhất, đơn giản một cách đánh lừa: liệu hướng tiếp cận học máy tăng cường bằng vệ tinh vốn báo cáo các giá trị R\textsuperscript{2} trên 0,8 cho Trung Quốc, Hoa Kỳ, và châu Âu có thể mang lại độ chính xác dự đoán không gian tương đương cho Việt Nam, với một mạng lưới thưa hơn nhiều so với những mạng lưới mà nó được phát triển trên đó? Các chương trước đã trả lời câu hỏi đó là không đối với con số nổi bật nhưng một cách chi tiết đáng kể đối với mọi thứ bên dưới nó. Chương 4 đã giải thích những gì mô hình học và vì sao nó hành xử như vậy; Chương 5 đã trình bày đánh giá định lượng đầy đủ từ kỳ vọng ban đầu đến bài toán trọng tâm. Chương kết luận này gắn kết các kết quả đó lại với nhau. Mục 6.1 tổng hợp các phát hiện của toàn bộ luận văn và gắn từng phát hiện trở lại bốn đóng góp đã nêu ở Chương 1. Mục 6.2 chuyển các phát hiện đó thành các khuyến nghị cụ thể cho chiến lược quan trắc chất lượng không khí của Việt Nam. Mục 6.3 tập hợp các hạn chế chính làm điều kiện cho các kết luận này, Mục 6.4 đề ra các hướng phát triển trong tương lai mà các kết quả thúc đẩy trực tiếp nhất, và một đoạn kết ngắn nêu lại sợi chỉ trọng tâm.

\section{Tổng hợp các phát hiện}

Kết quả trọng tâm của luận văn này là sự định lượng khoảng cách giữa các bài toán dự đoán, và nó là nội dung của đóng góp thứ nhất. Cùng một mô hình XGBoost-DART đạt R\textsuperscript{2} ngoài-fold khoảng 0,80 cho nội suy thời gian (đánh giá chéo ngẫu nhiên năm-fold) --- con số trực tiếp so sánh được với tài liệu khoa học đã công bố cho các khu vực được quan trắc dày đặc --- nhưng chỉ R\textsuperscript{2} trung bình theo trạm khoảng 0,20 cho ngoại suy không gian (đánh giá bỏ-một-trạm-ra), một khoảng cách khoảng 0,60. Khoảng cách này phản ánh sự khác biệt cơ bản giữa hai bài toán: nội suy thời gian được hưởng lợi từ việc biết lịch sử của trạm, trong khi ngoại suy không gian phải suy luận mức ô nhiễm nền của vị trí mà không có bất kỳ quan trắc nào tại đó. Đây là đánh giá có hệ thống đầu tiên đối chiếu kỳ vọng ban đầu với hiệu suất ngoại suy không gian thực tế cho PM2.5 tại Việt Nam, và khoảng cách này nhất quán với các phát hiện quốc tế: Kawano et al. \cite{kawano2025} thấy khoảng cách 0,18 với khoảng một nghìn trạm Ấn Độ, Meyer et al. \cite{meyer2018} thấy 0,66, và Ploton et al. \cite{ploton2020} thấy 0,39. Khoảng cách lớn hơn của Việt Nam phản ánh mạng lưới thưa hơn nhiều, và hệ quả tức thời là các con số nổi bật được báo cáo cho các mô hình PM2.5 ở Đông Nam Á và các mạng-lưới-thưa khác, khi chỉ đánh giá nội suy thời gian, không phản ánh năng lực ngoại suy không gian thực sự.

Phát hiện diễn giải quan trọng nhất, và trái tim của đóng góp thứ ba, là khả năng dự đoán ở Việt Nam tỷ lệ với mức ô nhiễm, và rằng kỹ thuật khai thác sự thật này kéo theo một phụ thuộc vòng tròn không thể triển khai. Việc phân tầng ba mươi bảy trạm huấn luyện thành bốn tầng theo PM2.5 trung bình của chúng và huấn luyện một mô hình riêng cho mỗi tầng --- sơ đồ T4F --- là đòn bẩy hiệu suất đơn lẻ lớn nhất được phát hiện trong nghiên cứu này, nâng trung bình theo trạm LOSO từ về cơ bản bằng không dưới một mô hình toàn cục mù-tầng lên con số nổi bật 0,20, một mức tăng khoảng +0,17 (trong-cùng-tệp, so với đường cơ sở không-phân-nhóm cùng-tệp 0,027) đến +0,26 --- có thể bảo vệ nhất là khoảng +0,17 --- vượt bất kỳ nhóm đặc trưng, lựa chọn siêu tham số, hay chiến lược tổ hợp nào được khảo sát. Phân tách theo tầng đơn điệu và rõ ràng: tầng ô nhiễm cao nhất t3 đạt R\textsuperscript{2} trung bình gần 0,56, tầng trung bình t2 gần 0,31, trong khi các tầng thấp và sạch t1 và t0 nằm tại hoặc dưới không, mô hình không thể vượt qua trung bình riêng của trạm ở nơi phương sai cần giải thích bản thân nó đã nhỏ. Việc sự phụ thuộc mức này là cơ chế vận hành, chứ không phải địa lý, chỉ có thể được đưa ra như một giả thuyết ở đây, bởi vì mức và vùng bị lẫn lộn ở đỉnh của phân bố: cả chín trạm t3 đều ở phía Bắc, nên gradient vùng biểu kiến và gradient mức không thể được tách bạch trong tầng cao nhất. Nơi duy nhất chúng có thể được so sánh ở mức cố định là t2, nơi hai vị trí Thành phố Hồ Chí Minh vượt điểm phần lớn các trạm phía Bắc có nồng độ tương tự và R\textsuperscript{2} trung bình t2 phía Nam (0,459, n=4) vượt phía Bắc (0,278, n=4); nhưng với chỉ bốn trạm mỗi bên, sự khác biệt không có ý nghĩa thống kê (Welch p=0,30, Mann-Whitney p=0,34, CI 95\% [-0,13,+0,49]), nên nó mang tính gợi ý chứ không phải xác nhận và không thể xác lập rằng vùng là không liên quan. Phụ thuộc là vòng tròn bởi vì việc gán một trạm vào tầng của nó đòi hỏi biết PM2.5 trung bình của nó, vốn chính xác là đại lượng mà mô hình tồn tại để dự đoán tại một vị trí không có quan trắc.

Cái giá của phụ thuộc vòng tròn đó định nghĩa ngưỡng trần triển khai, nửa sau của đóng góp thứ ba. Khi tầng oracle bị rút đi và chỉ các đặc trưng vệ tinh, khí tượng, và sử dụng đất thực sự quan sát được được phép, hiệu suất sụp đổ từ con số LOSO 0,20 xuống một trung vị theo trạm gần 0,04. Ba hướng tiếp cận triển khai được độc lập --- tổ hợp chuyên gia cổng-mềm, đường cơ sở dự đoán hai-pha, và dấu vân tay chế độ dẫn xuất từ vệ tinh --- hội tụ ở mép trên của một dải dương gần-không, bằng chứng mạnh rằng hiệu suất triển khai thực sự chỉ từ các biến quan sát được từ vệ tinh, khi không có hiệu chuẩn cục bộ, nhất quán với một dải dương gần-không mà mép trên của nó là khoảng +0,04 trung vị. Trên ngưỡng trần đó là oracle: việc tiêm trung bình thực của mỗi trạm giữ lại như một biên cơ sở cộng tính nâng trung bình theo trạm lên khoảng +0,27 với khoảng 92\% số trạm dương. Khoảng cách giữa dải triển khai được và ngưỡng trần oracle đó --- khoảng +0,28 đến +0,33 về R\textsuperscript{2} trung bình theo trạm và khoảng bốn mươi điểm phần trăm trạm dùng được --- là chi phí thông tin chính xác của việc không biết trung bình của một trạm. Bảy họ phương pháp riêng biệt được phát triển cụ thể để phá vỡ phụ thuộc --- ranh giới tầng mềm, các tương tác GHAP-lai, tổ hợp tầng, dự đoán tầng hai-pha, phân cụm dựa-trên-biến-quan-sát và định tuyến trạm-lân-cận-gần-nhất, loại bỏ đặc-trưng-bất-biến, và tái khung dị thường --- và cả bảy đều không đạt con số true-tier, mọi biến thể triển khai được rơi tại hoặc dưới nó. Các kết quả âm này không phải một suy nghĩ thêm; chúng tam giác đạc ranh giới của điều gì có thể đạt được và xác lập rằng không có đại diện quan sát được từ vệ tinh nào được kiểm tra giải quyết đáng tin cậy phụ thuộc. Do đó ngưỡng trần triển khai +0,04 là một ngưỡng trần thông tin cho \emph{ước tính quan sát được} về đường cơ sở, áp đặt bởi sự thưa thớt của mạng lưới mặt đất chứ không phải bởi bất kỳ sự thiếu sót nào về nỗ lực hay thuật toán.

Tuy nhiên, phụ thuộc quả thực thừa nhận một lời giải --- không phải từ các biến quan sát được, mà từ chính mạng lưới quan trắc. Bảy họ đều tìm trung bình của trạm trong không gian đặc trưng quan sát được; hướng tiếp cận thứ tám, định tuyến tiên nghiệm không gian, thay vào đó nội suy nó từ không gian vật lý, neo đường cơ sở của mỗi trạm giữ lại vào các trung bình quan trắc có-trọng-số-khoảng-cách của các trạm lân cận của nó và định tuyến các luồng dự đoán không-mục-tiêu về phía ước tính đó. Với GHAP được loại bỏ, quy trình hoàn toàn triển khai được này đạt R\textsuperscript{2} trung bình theo trạm 0,197 so với 0,203 của oracle --- khép gần như toàn bộ khoảng cách triển-khai-đến-oracle mà các đại diện quan sát được không thể --- và làm vậy với không có chuyển trạng thái cao-thành-không-cao và không có chuyển trạng thái thấp-thành-cao nguy hiểm nào, an toàn hơn một cách cận biên so với chính oracle. Trên bốn mươi sáu cảm biến chi phí thấp chưa thấy, nó giữ một trung bình theo trạm gần 0,20 mà không có vị trí sạch nào được ánh xạ thành ô nhiễm, tổng quát hóa như một bản đồ sàng lọc ô-nhiễm-cao. Điều này không mâu thuẫn với ngưỡng trần +0,04; nó đặc tả nó. Đường cơ sở không thể phục hồi từ các biến quan sát được từ vệ tinh và có thể phục hồi bằng nội suy không gian, nên kỹ năng triển khai có thể đạt được không phải một con số duy nhất mà là một hàm của sự gần kề với mạng lưới: gần-oracle ở nơi mục tiêu có các trạm lân cận khả dụng, sụp đổ về phía sàn quan sát được ở nơi không có. Ràng buộc then chốt, một lần nữa, là tầm với của mạng lưới mặt đất. Mô hình vẫn là một quy trình nghiên cứu nặng, giàu-đặc-trưng chứ không phải một bản đồ vận hành --- nó không được thiết kế cho suy luận phủ-liên-tục trực tiếp --- nên đóng góp là \emph{nguyên lý lời giải} đã được chứng minh, không phải một sản phẩm được triển khai.

Việc đánh giá có hệ thống các sản phẩm PM2.5 toàn cầu --- đóng góp thứ hai --- cung cấp kết quả âm bổ sung và giải thích vì sao các sản phẩm đó không thể cung cấp thông tin bị thiếu. Cả ba sản phẩm tái phân tích và mô hình vận chuyển hóa học đều thất bại như các bộ ước tính độc lập trên Việt Nam, nhưng mỗi sản phẩm thất bại theo một phương diện khác biệt đặc trưng. GEOS-CF thể hiện một chu kỳ ngày đêm đảo ngược so với quan trắc mặt đất, đạt cực tiểu vào giờ rạng đông khi PM2.5 bề mặt đạt đỉnh, cùng với một sai số trung bình theo trạm vượt +244\% --- chữ ký không thể nhầm lẫn của một kiểm kê phát thải và sơ đồ lớp biên được tinh chỉnh cho một mẫu ngày đêm khác. MERRA-2 xấp xỉ cân bằng khối lượng về trung bình nhưng trơ về thời gian, với một Chỉ số Phù hợp theo trạm gần 0,39 đến 0,42, nhất quán với vị thế đã công bố của nó là kém nhất trong bất kỳ khu vực toàn cầu nào. GHAP mang gradient không gian tốt nhất trong ba, gán đúng tầng ô nhiễm của chỉ khoảng một nửa số trạm, nhưng đặt trên một sàn bị nâng cao tại các vị trí nông thôn sạch, ước tính vượt các nồng độ sạch nhất khoảng +10 $\mu$g/m\textsuperscript{3} và tích cực sắp sai thứ tự chúng. Được đánh giá như các bản đồ tĩnh, cả GHAP và MERRA-2 đều chia sẻ cùng bệnh lý vị-trí-sạch này --- đúng chế độ thất bại mà mô hình của chính luận văn thể hiện và vì cùng một lý do cấu trúc: ở nơi tín hiệu bề mặt thực sự nhỏ so với nhiễu truy xuất và tái phân tích, các sản phẩm được xây dựng chủ yếu từ AOD vệ tinh kế thừa một sàn dương và mất khả năng phân biệt giữa các nồng độ thấp. Sợi chỉ chung là cả ba đều được hiệu chuẩn chủ yếu trên dữ liệu Đông Á và phương Tây dày đặc và không chuyển giao một cách sạch sẽ đến môi trường phát thải nhiệt đới, gió mùa, xe-máy-và-sinh-khối của Việt Nam. Không sản phẩm nào có thể được coi là một tiên nghiệm đáng tin cậy cho PM2.5 theo giờ mà không có hiệu chỉnh cục bộ, đó chính xác là lý do mô hình hoàn chỉnh dựa vào các trạm lân cận mặt đất và các quan trắc vệ tinh trực tiếp thay vì đầu ra CTM, và vì sao các đặc trưng CTM đóng góp về cơ bản không độ tăng nào, việc loại bỏ chúng là trung tính hoặc có lợi.

Đánh giá độc lập --- đóng góp thứ tư --- cung cấp đối ứng dương cho các kết quả âm này và chỉ ra ràng buộc then chốt thực sự. Nó tốt nhất được hiểu là một phép kiểm tra chuyển đặc trưng từ trạm gần nhất cộng với sự phù hợp cảm biến chứ không phải dự đoán không gian thực sự tại các vị trí chưa thấy: đối với mỗi cảm biến chi phí thấp giữ lại, các đặc trưng vệ tinh, AOD, TROPOMI, và khí tượng ERA5 được lấy từ bản ghi theo giờ của trạm quan trắc (KK) gần nhất, và chỉ có PM2.5 mục tiêu đến từ chính cảm biến, trong khi các điểm neo không gian RFSI được định vị đúng. Được kiểm tra theo cách này so với ba mươi chín cảm biến chi phí thấp độc lập và trạm tham chiếu cấp-chuẩn Đại sứ quán Hoa Kỳ, không trạm nào được dùng trong phát triển mô hình, mô hình vùng đồng bằng sông Hồng đạt một R\textsuperscript{2} trung vị 0,53 trên các cảm biến chi phí thấp với khoảng 85\% số vị trí dương, và R\textsuperscript{2} = 0,68 tại thiết bị suy giảm beta của Đại sứ quán Hoa Kỳ --- kết quả độc lập đáng tin cậy nhất duy nhất trong luận văn, mặc dù nó một phần là kết quả chuyển đặc trưng và nên được đọc như vậy. Các con số này cạnh tranh với các kết quả đánh giá chéo không gian tốt nhất được công bố trên quốc tế và khép khoảng cách với mốc chuẩn Ấn Độ trong vùng phủ dày đặc, nơi R\textsuperscript{2} trung vị của các vị trí trong vòng 10 km đạt 0,64. Khoảng cách rộng giữa trung vị mạnh và trung bình yếu là chữ ký giờ-đã-quen của một số ít vị trí xa, bị suy giảm, và nó bản thân là câu chuyện mật độ: kỹ năng dự đoán suy giảm rõ ràng theo khoảng cách đến trạm neo gần nhất (mối quan hệ suy-giảm-theo-khoảng-cách được dẫn xuất ở Mục 4.5 và 5.3), một sự suy giảm phần nào phản ánh sai số thay-thế-đặc-trưng tăng dần phát sinh khi khí tượng của trạm quan trắc gần nhất thay thế cho các điều kiện ở xa hơn. Tầm quan trọng đặc trưng xác nhận cơ chế --- nội suy không gian từ các trạm lân cận mặt đất (RFSI) chiếm tỷ lệ lớn nhất của độ tăng mô hình ở khoảng 37\%, quan trắc vệ tinh đứng sát thứ hai gần 31\%, khí tượng khoảng 22\%, các mã hóa thời gian gần 10\%, và đầu ra CTM về cơ bản không gì cả --- nên biến dự đoán mạnh nhất của mô hình, theo thiết kế, là PM2.5 của các trạm lân cận. Do đó kết luận thống nhất trên cả bốn đóng góp là duy nhất và định lượng: kỹ năng dự đoán không gian ở Việt Nam bị chặn trên bởi tầm với của mạng lưới mặt đất chứ không phải bởi thuật toán hay vệ tinh. Năng lực không gian của mô hình bị giới hạn bởi mật độ trạm, và mật độ đó có thể được mở rộng.

\section{Khuyến nghị cho chiến lược quan trắc chất lượng không khí của Việt Nam}

Các phát hiện của luận văn này chuyển thành bốn khuyến nghị cụ thể về cách Việt Nam nên tạo ra, kiểm chứng, và truyền đạt các ước tính PM2.5 dựa trên vệ tinh.

Thứ nhất, phương pháp đánh giá phải phù hợp với bài toán đang được giải quyết. Khi mục đích là triển khai mô hình tại các vị trí không có quan trắc --- tức bài toán ngoại suy không gian --- thì đánh giá bỏ-một-trạm-ra hoặc đánh giá chéo không gian tương đương là phương pháp phù hợp, không phải đánh giá chéo ngẫu nhiên. Đánh giá chéo ngẫu nhiên đo nội suy thời gian (R\textsuperscript{2} khoảng 0,80), trong khi ngoại suy không gian chỉ đạt R\textsuperscript{2} khoảng 0,20 --- khoảng cách khoảng 0,60 không phải là lỗi phương pháp mà phản ánh hai bài toán khác nhau về bản chất. Nhưng khi một R\textsuperscript{2} từ bài toán dễ hơn (nội suy thời gian) được dùng để biện minh cho năng lực trong bài toán khó hơn (ngoại suy không gian), nó gây hiểu lầm một cách hệ thống. Đối với một đất nước được quan trắc thưa thớt như Việt Nam, chỉ con số ngoại suy không gian mới có ý nghĩa vận hành. Các đặc tả mua sắm, tiêu chí chấp nhận quy định, và các khẳng định học thuật về độ chính xác không gian đều nên đòi hỏi mô hình được huấn luyện lại với trạm đánh giá bị loại khỏi cả tập huấn luyện lẫn tập trạm-lân-cận-không-gian, và nên báo cáo trung vị theo trạm và tỷ lệ trạm có R\textsuperscript{2} dương cùng với bất kỳ con số gộp nào.

Thứ hai, các sản phẩm PM2.5 mô hình vận chuyển hóa học và vệ tinh toàn cầu --- trong đó có GEOS-CF, MERRA-2, và GHAP --- không được dùng như các tiên nghiệm chưa hiệu chỉnh cho Việt Nam, dù là các đặc trưng đầu vào được đặt trọng số một cách tin tưởng, các bề mặt phơi nhiễm cho đánh giá sức khỏe, hay các bộ lấp khoảng trống ở nơi không có dữ liệu mặt đất. Đánh giá có hệ thống trong luận văn này cho thấy mỗi sản phẩm thất bại theo một phương diện riêng, và các thất bại nghiêm trọng nhất chính xác tại các vị trí nông thôn sạch nơi, vì thiếu trạm mặt đất, một nhà hoạch định sẽ muốn tin một bản đồ toàn cầu nhất. Một sản phẩm thổi phồng một mức 10 $\mu$g/m\textsuperscript{3} thực thành 20 $\mu$g/m\textsuperscript{3} và sắp sai thứ tự các vị trí sạch nhất không chỉ thêm nhiễu; nó tạo ra sai số tự tin chính xác ở nơi nó được hỗ trợ ít nhất. Ở nơi các sản phẩm này được dùng, chúng chỉ nên được đưa vào sau khi hiệu chỉnh độ thiên lệch cục bộ tường minh đối chiếu với các phép đo mặt đất Việt Nam, và các giá trị tuyệt đối của chúng --- đặc biệt chu kỳ ngày đêm đảo ngược và độ thiên lệch cao gấp nhiều lần của GEOS-CF --- không bao giờ nên được báo cáo như các ước tính PM2.5 mà không có hiệu chỉnh đó.

Thứ ba, và quan trọng nhất về hệ quả, Việt Nam nên theo đuổi một mạng lưới quan trắc lai gồm các trạm chuẩn-tham-chiếu thưa thớt được làm dày bằng các cảm biến chi phí thấp đã hiệu chuẩn, với mục tiêu hoạch định tường minh rằng không khu vực có dân cư nào nằm cách một trạm neo quá khoảng 10 km. Bằng chứng suy-giảm-theo-khoảng-cách rõ ràng: trong vòng 10 km của một trạm huấn luyện mô hình cạnh tranh với các kết quả quốc tế tốt nhất, và ngoài 15 km các dự đoán của nó suy giảm về phía vô dụng. Các cảm biến chi phí thấp đi vào hoạt động cuối năm 2025 chỉ đóng vai trò là các mục tiêu đánh giá độc lập trong luận văn này, không phải các điểm neo nội suy, nên luận điểm làm dày dựa trên phép chiếu chứ không phải một mức tăng đã được chứng minh: ngoại suy đường cong suy-giảm-theo-khoảng-cách đo được, việc mở rộng độ phủ trạm neo sao cho phần lớn đất nước có dân cư rơi vào chế độ 10 km sẽ, theo phép chiếu đó, đưa phần lớn Việt Nam vào dải nơi mô hình đã đạt một R\textsuperscript{2} trung vị trên 0,5. Việc thêm các cảm biến chi phí thấp đã hiệu chuẩn vào tập trạm-lân-cận-không-gian là cách tự nhiên để hiện thực hóa độ phủ như vậy, nhưng độ lớn của mức tăng thu được vẫn còn phải được chứng minh thực nghiệm. Đây là một can thiệp khả thi và rẻ tiền hơn nhiều so với việc nhân lên mạng lưới quan trắc. Các trạm chuẩn-tham-chiếu vẫn thiết yếu như các điểm neo hiệu chuẩn và kiểm soát chất lượng để các cảm biến chi phí thấp được hiệu chỉnh đối chiếu, nhưng chính tầm với của mạng lưới được làm dày, chứ không phải số lượng thiết bị quan trắc, mới chi phối độ chính xác có thể đạt được.

Thứ tư, các sản phẩm vận hành dẫn xuất từ một mạng lưới như vậy nên được công bố như các bản đồ nhận-biết-độ-tin-cậy với độ bất định biến thiên không gian chứ không phải như các giá trị xác định duy nhất. Luận văn đã chứng minh rằng cả mô hình và các sản phẩm toàn cầu hiện có đều tự tin sai nhất chính xác ở nơi chúng được hỗ trợ ít nhất --- xa các trạm mặt đất và ở nồng độ thấp --- và rằng một bản đồ đơn-giá-trị không thể truyền đạt điều này. Thay vào đó mỗi ô lưới nên mang một độ bất định tường minh mở rộng theo khoảng cách đến trạm neo gần nhất và ở nồng độ thấp, sao cho một dự đoán trong một khu vực ô nhiễm được neo tốt được trình bày như đáng tin cậy trong khi một dự đoán trong một khu vực sạch cô lập được gắn cờ như bất định. Sự tái khung này chuyển các hạn chế đã được ghi nhận của phương pháp từ một trách nhiệm pháp lý thành thông tin trung thực, có thể hành động cho các người dùng y tế công cộng và quy định phụ thuộc vào các bản đồ, và nó theo trực tiếp từ ràng buộc mật độ chi phối kỹ năng nền tảng.

\section{Hạn chế}

Một số hạn chế làm điều kiện cho các kết luận ở trên và được tập hợp ở đây bằng cách tham chiếu chứ không phải dẫn xuất lại. Đánh giá độc lập cảm-biến-chi-phí-thấp là kết quả được lưu ý nhiều nhất: như đã nêu ở Mục 6.1, nó là một phép kiểm tra chuyển đặc trưng từ trạm gần nhất cộng với sự phù hợp cảm biến chứ không phải dự đoán không gian thực sự tại các vị trí chưa thấy, vì mọi đặc trưng vệ tinh, AOD, TROPOMI, và khí tượng ERA5 cho một cảm biến giữ lại được thay thế từ bản ghi của trạm quan trắc gần nhất trong khi chỉ có PM2.5 mục tiêu là độc lập; do đó một phần của cả R\textsuperscript{2} độc lập nổi bật lẫn sự suy-giảm-theo-khoảng-cách mà nó thể hiện đều quy cho sự thay thế đặc trưng. Đánh giá đó cũng bị giới hạn ở mùa khô, để lại một nhiễu loạn theo mùa chưa được định lượng, và dựa trên các bản ghi chi phí thấp tán xạ laser ngắn mà sự ổn định của chúng qua các giai đoạn dài hơn chưa được kiểm tra. Trung bình độc lập mạnh còn bị dẫn dắt bởi một số ít vị trí và nhạy với một giá trị ngoại lai duy nhất, nên trung vị là bản tóm tắt đáng tin cậy hơn, và các con số nổi bật được tính trên tập con vị trí đầy-đủ-đặc-trưng mà tất cả các đầu vào sẵn có chứ không phải trên toàn bộ tập cảm biến.

Các hạn chế còn lại liên quan đến mạng lưới mặt đất và các quy ước báo cáo. Mạng lưới quan trắc nhỏ --- ba mươi bảy trạm trên toàn quốc, mười hai trong mô hình đồng bằng --- vốn vừa hạn chế tầm với không gian có thể đạt được vừa giới hạn sức mạnh thống kê của mọi so sánh theo-tầng và theo-vùng, như sự tương phản Bắc--Nam t2 không có ý nghĩa ở Mục 6.1 minh họa. Xuyên suốt luận văn, các bản tóm tắt theo trạm được báo cáo như các trung bình và trung vị có-trọng-số-n của R\textsuperscript{2} theo trạm chứ không phải như một R\textsuperscript{2} gộp duy nhất được tính trên tất cả các giờ cùng một lúc; như đã lưu ý ở Mục 5.1, hai quy ước này không nhất thiết phải thống nhất, và các con số theo trạm là các con số có ý nghĩa vận hành nhưng không thể hoán đổi với một giá trị gộp. Các hạn chế này được thừa nhận một cách thẳng thắn: không hạn chế nào lật ngược phát hiện khoảng-cách-đánh-giá trọng tâm, nhưng mỗi hạn chế giới hạn việc các con số độc lập và vùng thuận lợi hơn nên được tổng quát hóa xa đến đâu.

\section{Hướng phát triển trong tương lai}

Các kết quả của luận văn này chỉ ra một số hướng nghiên cứu cụ thể, được tổ chức quanh phát hiện trọng tâm rằng ngưỡng trần triển khai được đặt bởi chi phí thông tin của việc không biết PM2.5 trung bình của một trạm.

Hướng công việc trực tiếp nhất là mở rộng bộ ước tính đường-cơ-sở-không-gian ra ngoài tầm với của mạng lưới hiện tại. Luận văn này đã thực hiện bước quyết định mà dàn ý trước đó chỉ dự đoán như công việc tương lai: quy trình định tuyến tiên nghiệm không gian (Mục 3.6.4 và 5.5) ước tính đường cơ sở của một trạm bằng cách nội suy các trung bình quan trắc của các trạm lân cận của nó và, được dùng để neo và định tuyến các luồng không-mục-tiêu, khép gần như toàn bộ khoảng cách triển-khai-đến-oracle (trung bình theo trạm 0,197 so với 0,203 của oracle). Hạn chế của bộ ước tính đó nay được định nghĩa sắc nét: nó hoạt động ở nơi mục tiêu có các trạm lân cận khả dụng và suy giảm ở nơi không có, nên bài toán mở không còn là phục hồi đường cơ sở gần mạng lưới mà là ước tính nó ở nơi mạng lưới vắng mặt. Hai lộ trình bổ sung theo sau. Lộ trình thứ nhất là một bộ ước tính trung-bình-trạm quan sát được giàu hơn --- kết hợp phân loại sử dụng đất tinh hơn, mật độ dân số và môi-trường-xây-dựng theo lưới, các đại diện mạng-lưới-đường và phát-thải-giao-thông, và các kiểm kê phát thải từ-dưới-lên thay vì các khí hậu vệ tinh thô được kiểm tra ở đây --- để nâng ước tính đường cơ sở ở các vùng mà tiên nghiệm không gian không thể vươn tới; hồi quy vệ tinh hai-pha của luận văn này chỉ đạt một R\textsuperscript{2} bỏ-một-ra khoảng 0,59, và việc cải thiện nó là con đường mở rộng độ phủ vào các khu vực được quan trắc thưa thớt. Lộ trình thứ hai là tái thiết kế bản thân quy trình: mô hình hiện tại là một hiện vật nghiên cứu nặng, đa-luồng, giàu-đặc-trưng, không phải một cỗ máy suy luận, và việc biến nguyên lý lời giải đã được chứng minh thành một bản đồ vận hành phủ-liên-tục sẽ đòi hỏi chưng cất nó thành một bộ ước tính triển khai được duy nhất đủ nhẹ để chạy trên một lưới quốc gia.

Một lộ trình bổ sung và có khả năng mạnh hơn là hiệu chuẩn tại-chỗ vài-mẫu: thu thập một bản ghi mặt đất ngắn --- cỡ vài tuần --- tại một vị trí mục tiêu để ước tính PM2.5 trung bình của nó trực tiếp, rồi cung cấp trung bình thực nghiệm đó cho mô hình như một biên cơ sở. Luận văn này xác lập rằng việc biết trung bình trạm thực nâng hiệu suất lên mức oracle khoảng +0,27, và một triển khai ngắn của một cảm biến chi phí thấp đã hiệu chuẩn là cách rẻ nhất để thu được một ước tính dùng được về trung bình đó mà không cần một trạm thường trực. Việc định lượng một bản ghi như vậy có thể ngắn đến đâu trước khi trung bình ước tính làm suy giảm hiệu suất --- và độ chính xác thu được nội suy giữa sàn triển khai +0,04 và ngưỡng trần oracle +0,27 như một hàm của độ dài bản ghi ra sao --- là một thí nghiệm được định nghĩa rõ và có giá trị cao.

Hướng thứ ba là chính thức hóa bản đồ độ tin cậy hoặc tin tưởng vệ tinh trên một mô hình toàn cục duy nhất. Luận văn cho thấy một mô hình triển khai được hoạt động tốt hơn nhiều tại tầng ô nhiễm so với tại các vị trí sạch, và rằng các biến quan sát được từ vệ tinh có thể phân loại chế độ của một trạm một cách không hoàn hảo nhưng giàu thông tin. Thay vì định tuyến các dự đoán qua các mô hình riêng theo tầng, một mô hình toàn cục duy nhất có thể phát ra một điểm số độ tin cậy theo vị trí dẫn xuất từ phân loại chế độ quan-sát-được-từ-vệ-tinh của nó, trực tiếp vận hành hóa việc lập bản đồ nhận-biết-độ-tin-cậy được khuyến nghị ở Mục 6.2 và cung cấp độ bất định biến thiên không gian mà các bản đồ đơn-giá-trị thiếu.

Hai hướng đi nữa theo sau từ cấu trúc mạng lưới và dữ liệu. Các mô hình vùng cho các cụm dày đặc nên được phát triển ngoài đồng bằng sông Hồng ở bất cứ đâu mật độ trạm cho phép: trung bình LOSO 0,302 và trung vị 0,433 của mô hình đồng bằng vượt một cách khiêm tốn cấu hình toàn quốc triển khai được mù-tầng trên cùng các trạm, ủng hộ chiến lược vùng và gợi ý rằng bất kỳ cụm tương lai đủ dày nào --- quanh Thành phố Hồ Chí Minh khi mạng lưới chi phí thấp mở rộng về phía Nam, chẳng hạn --- sẽ hưởng lợi từ một mô hình chuyên dụng chia sẻ khí tượng, các loại phát thải, và phơi nhiễm xuyên biên giới. Cuối cùng, đánh giá cảm-biến-chi-phí-thấp trong luận văn này bị giới hạn ở mùa khô, nên các kết quả độc lập mạnh mang một nhiễu loạn theo mùa chưa được định lượng; việc thu thập các bản ghi chi phí thấp dài hơn, đa-mùa sẽ cho phép các con số mùa khô được tách khỏi kỹ năng không gian quanh năm thực sự và sẽ loại bỏ lưu ý quan trọng nhất gắn với đánh giá độc lập.

\section{Lời kết}

Luận văn này bắt đầu với khoảng cách giữa một con số được công bố và một con số trung thực, và nó kết thúc bằng việc đã định vị nguồn gốc của khoảng cách đó một cách chính xác. R\textsuperscript{2} khoảng 0,80 mà mô hình tăng cường bằng vệ tinh đạt được dưới đánh giá chéo ngẫu nhiên là thực nhưng đo sai bài toán; dưới đánh giá bỏ-một-trạm-ra trung thực con số là khoảng 0,20, và một khi phụ thuộc vòng tròn vào mức ô nhiễm riêng của một trạm bị rút đi, ước tính thực sự triển khai được chỉ từ các biến quan sát được từ vệ tinh là một trung vị theo trạm gần 0,04. Các sản phẩm mô hình vận chuyển hóa học và vệ tinh vốn có thể đã cung cấp thông tin bị thiếu thất bại trên Việt Nam theo ba cách riêng biệt, và bảy họ phương pháp đại diện quan sát được không thể phục hồi nó --- nhưng đại lượng bị thiếu hóa ra có thể phục hồi từ mạng lưới quan trắc chứ không phải từ các biến quan sát được: việc neo đường cơ sở vào các trạm lân cận thông qua định tuyến tiên nghiệm không gian nâng trung bình theo trạm triển khai được trở lại 0,197, chỉ cách một sợi tóc so với 0,203 của oracle, và mang đến các cảm biến chi phí thấp chưa thấy như một bản đồ sàng lọc an toàn. Trong vùng phủ mặt đất dày đặc, mô hình đạt một R\textsuperscript{2} trung vị 0,53 và một R\textsuperscript{2} cấp-tham-chiếu 0,68, cạnh tranh với các kết quả quốc tế tốt nhất. Bài học thống nhất là kỹ năng dự đoán không gian cho PM2.5 ở Việt Nam bị giới hạn không phải bởi thuật toán và không phải bởi vệ tinh mà bởi tầm với của mạng lưới mặt đất --- phụ thuộc bị phá vỡ chính xác ở nơi mạng lưới vươn tới và ràng buộc chính xác ở nơi không --- một ràng buộc mà, đáng khích lệ, là ràng buộc khả thi nhất để nới lỏng. Việc làm dày mạng lưới đó bằng các cảm biến chi phí thấp đã hiệu chuẩn sao cho không khu vực có dân cư nào nằm xa một trạm neo, đánh giá một cách trung thực, và lập bản đồ với độ tin cậy tường minh cùng nhau cấu thành một con đường thực tiễn và phải chăng từ sàn triển khai được ghi nhận ở đây hướng tới độ chính xác mà quan trắc chất lượng không khí của Việt Nam sẽ đòi hỏi.
